% Figure: family of polytropic processes p V^n = const leaving a common start
% state, for n = 0 (isobaric), 1 (isothermal), gamma=1.4 (adiabatic), and
% n large (toward isochoric). Drawn with pgfplots.
\begin{figure}[htbp]
\centering
\begin{tikzpicture}
\begin{axis}[
    width=0.78\textwidth, height=0.52\textwidth,
    xlabel={volume $V$ (arb.\ units)},
    ylabel={pressure $p$ (arb.\ units)},
    domain=1:4, samples=80,
    xmin=0.8, xmax=4.4, ymin=0, ymax=4.6,
    legend pos=north east, legend cell align=left,
    legend style={font=\scriptsize, draw=enggray},
    tick label style={font=\scriptsize},
    label style={font=\small},
    axis lines=left,
]
  % all curves leave the common start (V=1, p=4)
  % isobaric n=0: p constant = 4
  \addplot[engamber, thick, domain=1:4] {4};
  \addlegendentry{$n=0$ isobaric}
  % isothermal n=1: p = 4/V
  \addplot[engblue, thick, domain=1:4] {4/x};
  \addlegendentry{$n=1$ isothermal}
  % adiabatic n=gamma=1.4: p = 4 * V^(-1.4)
  \addplot[engred, thick, domain=1:4] {4*x^(-1.4)};
  \addlegendentry{$n=1.4$ adiabatic}
  % polytropic n=2: p = 4 * V^(-2), a real intercooled-ish compression
  \addplot[englime!70!engblack, thick, domain=1:4] {4*x^(-2)};
  \addlegendentry{$n=2$ polytropic}
  % common start point marker
  \addplot[only marks, mark=*, engblack, mark size=2pt] coordinates {(1,4)};
\end{axis}
\end{tikzpicture}
\caption[Family of polytropic processes on pressure-volume axes]{Polytropic
processes $pV^n = \text{const}$ leaving one common start state on
pressure-volume axes. The exponent $n$ slides between the named limits:
$n=0$ is the horizontal isobar, $n=1$ is the isotherm, $n=\gamma=1.4$ is the
reversible adiabat, and $n\to\infty$ would be the vertical isochore. A real
compression with some heat loss to the cylinder walls sits between the
adiabat and the isotherm; the case $n=2$ shown here is steeper than the
adiabat, the path of a compression that loses heat faster than an adiabatic
one. The area under each curve is the boundary work, so reading off the
ordering of the areas ranks the work of the four processes for the same
volume change.}
\label{fig:vol04ch01:polytropic}
\end{figure}
